\section{Problem Formulation}
\label{sec:problem}

\subsection{Agent Model}

We model a coding agent as a function $\mathcal{A}$ that maps a context
$c = (s, h)$---a system prompt $s$ and conversation history $h$---to a response
$r \in \mathcal{R}$, where $\mathcal{R}$ includes text (code, explanations) and
structured actions (tool calls, file writes). The agent interacts with a user
through a sequence of turns.

A \emph{clarification event} occurs when the agent emits a response of the form
$r = \texttt{CLARIFYING\_QUESTION}: q$ for some question $q$, and the subsequent
user turn $u$ (the clarification answer) is appended to $h$. The agent then
produces a final response based on the augmented context.

\subsection{Attack Model}

An adversary $\mathcal{V}$ has write access to exactly one input channel. We
consider two channels:

\begin{description}
  \item[Tool output ($\mathcal{C}_{\mathrm{exec}}$)] The adversary injects content
    into a simulated tool output (e.g., the result of a file read). This represents
    the standard prompt-injection threat model for tool-using agents.
  \item[Clarification answer ($\mathcal{C}_{\mathrm{clar}}$)] The adversary injects
    content into the user's clarification answer $u$. The payload is appended to
    (or replaces part of) the legitimate answer.
\end{description}

In both cases, the injected payload $p$ contains an instruction that, if followed,
produces a mechanically detectable harmful behavior in the agent's output.

\subsection{The Clarification Tax}

For a fixed payload $p$ and scenario $\sigma$, let $\asr(\mathcal{A}, p, \mathcal{C})$
denote the probability that the agent $\mathcal{A}$ follows payload $p$ when it
arrives through channel $\mathcal{C}$. We define the \emph{clarification tax} as:
\begin{equation}
  \tax(\mathcal{A}, p) = \asr(\mathcal{A}, p, \mathcal{C}_{\mathrm{clar}})
                       - \asr(\mathcal{A}, p, \mathcal{C}_{\mathrm{exec}})
  \label{eq:tax}
\end{equation}

A positive $\tax$ means the agent is more susceptible to injection through the
clarification channel. The ASPI report~\citep{aspi2026} provides evidence that
$\tax > 0$ for general agents; we test whether this holds for coding agents and
whether it can be reduced.

\subsection{Hypotheses}

We test four hypotheses:

\begin{description}
  \item[H1 (Replication)] $\tax > 0$ for coding agents across model families.
  \item[H2 (Channel vs.~State)] The tax is partly attributable to the agent's
    \emph{interaction state} (having just asked a question), not only to the delivery
    channel.
  \item[H3 (Mitigation)] Provenance-aware defenses reduce $\tax$ significantly.
  \item[H4 (Trade-off)] Stronger defenses impose greater collaboration cost
    (measured by $\ccr$, the clarification compliance rate).
\end{description}

\paragraph{Metrics.} We report:
\begin{itemize}
  \item $\asr$ -- Attack Success Rate: fraction of trials where the injected behavior
    appears in the agent's output.
  \item $\tsr$ -- Task Success Rate: fraction of trials where the benign task is
    correctly completed.
  \item $\ccr$ -- Clarification Compliance Rate: fraction of trials where the benign
    part of the clarification answer is correctly incorporated.
  \item $\bar{T}$ -- Average turns to clarification.
  \item $\tax$ -- Clarification tax (Equation~\ref{eq:tax}).
\end{itemize}
All estimates include 95\% bootstrap confidence intervals over $N=3$ runs per
scenario per condition.
